% Results — honest synthesis (all numbers from protocols/runs/, 2026-08-23)
\documentclass[preview]{standalone}
\begin{document}
\section{Results}

\subsection{Protocol machinery works end-to-end}
The full pipeline (guarded rollout identity, evidence bundles, prequential
streams, three-channel cost ledger, SafeCommit gate) ran on CTS, AppWorld
and tau2 with Qwen3-4B and Mistral-7B. All correctness gates passed
(R001-R003; Guard contract 24/24).

\subsection{SafeCommit reduces catastrophic updates}
In candidate-archive replay over benign/mixed/poisoned/abrupt-shift
streams, the empirical-Bernstein e-process gate commits zero harmful
candidates (catastrophic-update rate relative reduction 100\% vs
always-commit) while keeping benign/mixed commit rates at 0.11/0.10
(non-degenerate).

\subsection{Updates do not produce a stable prequential gain}
\begin{tabular}{lcccc}
\toprule
Setup & frozen & naive & egc & n \\
\midrule
CTS 8-task (Qwen3-4B) & 0.625 & 0.500 & 0.500 & 2-3 seeds \\
AppWorld 3-task (Qwen3-4B) & 0.000 & 0.000 & --- & 1 \\
tau2 8-task (Mistral-7B) & 0.072 & 0.093 & 0.077 & 2 \\
tau2 16-task (Mistral-7B) & 0.025 & 0.013 & --- & 2 \\
\bottomrule
\end{tabular}

LoRA-GRPO updates (4 steps, lr 5e-6, per-task batches) do not improve
first-attempt performance on future tasks under matched budgets; effect
directions flip with task pool. Single-step updates produced zero logit
drift (R002 canary); the 4-step variant changes behavior measurably but
not in a direction that transfers.

\subsection{Interpretation}
The paper contributes (a) a validated protocol machinery for
deployment-period agent RL under partial evidence (evidence tiers,
prequential induction, budget accounting, risk-controlled commit), and
(b) a reproducible negative result: at these scales and models,
test-time LoRA-RL does not yield stable inductive transfer, while
risk-controlled commit eliminates catastrophic updates. The harness is
open for stronger backbones and longer streams.
\end{document}
